%-------------------------------------------------------------------------------
% SECTION TITLE
%-------------------------------------------------------------------------------
\cvsection{Actividades Extracurriculares}
%-------------------------------------------------------------------------------
% CONTENT
%-------------------------------------------------------------------------------
\begin{cventries}

%\cventry
%{Estudiante}
%{Udemy, (El curso completo de desarrollo web)}
%{Barcelona, España} % Location
%{Jun. 2021 - Ago. 2021}
%{
%      \begin{cvitems} % Description(s) of experience/contributions/knowledge
%        \item {Obtuve experiencia enfocada en aplicaciones y software web.}
%        \item {Practiqué JavaScript, HTML y CSS.}
%      \end{cvitems}
%}

% \cventry
% {Ingeniero de Desarrollo de Software}
% {Hackató #bitsxlaMarató de TV3 2022}
% {Ene 2022}
% {Barcelona, España}
% {
% \begin{cvitems}
% \item{Desarrollé un detector de aneurismas basado en una red neuronal convolucional utilizando MaskRCNN, que implicó la creación de un modelo 3D de la aorta de una rata y el cálculo de la sección más ancha para detectar un aneurisma.}
% \item{Gané la hackathon #bitsxlaMarató de TV3 2022 "per la salut cardiovascular" por este proyecto.}
% \item{Diseñé e implementé una interfaz fácil de usar para cargar cualquier video y detectar la aorta, así como para visualizar el modelo 3D.}
% \item{Utilicé habilidades en desarrollo de software y visión por computador para lograr los objetivos del proyecto.}
% \end{cvitems}
% }

\cventry
{Ingeniero de Desarrollo de Software}
{Hackató \#bitsxlaMarató de TV3 2022}
{Barcelona, España} % Location 
{Ene 2022}
{
\begin{cvitems}
\item{Desarrollé un detector de aneurismas basado en una red neuronal convolucional utilizando MaskRCNN para la salud cardiovascular.}
\item{Gané el hackathon \#bitsxlaMarató de TV3 2022 por este proyecto.}
\item{Diseñé e implementé una interfaz fácil de usar para cargar cualquier video y detectar la aorta, así como para visualizar el modelo 3D.}
\end{cvitems}
}


\cventry
{Contribuidor de la comunidad de código abierto}
{Automatización y documentación de configuraciones de Linux}
{Barcelona, España} % Location
{Jun. 2017 - Presente}
{
      \begin{cvitems} % Description(s) of experience/contributions/knowledge
        \item {Me especialicé en scripting de shell y automatización de múltiples tareas a través de herramientas como Stow o Rclone.}
        \item {Aprendí cómo documentar soluciones a problemas encontrados mediante markdown/LaTeX.}
        \item {Los sistemas de control de fuente como Git se volvieron indispensables para mi flujo de trabajo a medida que las tareas se volvían más grandes.}
      \end{cvitems}
}
\end{cventries}